\subsection{Namelist Input}
\label{SEC:Files:Namelist_Input}

GFR uses a Fortran namelist file for specifying user inputs.
%
GFR is executed on the command line with the following syntax:
%
\begin{InputExample}
  \item[{\inptext{[MPI commands] gfr [optional args]}}]
\end{InputExample}
%
where \inptext{[MPI~commands]} specifies the MPI executable and MPI arguments if
running in parallel, \inptext{gfr} is the GFR executable, and 
\inptext{[optional~args]} are the optional arguments that will be processed by
GFR.
%
The logic for detecting and verifying the GFR input file is as follows:
%
\begin{enumerate}
  \item
    If the number of optional arguments is 1, it tests to see if the argument
    represents an existing file containing GFR namelists.
    %
  \item
    If the number of optional arguments is greater than 1, it searches through
    the arguments to find one equal to \quoterm{\inptext{-f}}.
    %
    If the input flag \quoterm{\inptext{-f}} is found as one of the arguments,
    it tests to see if the very next argument following the flag represents an
    existing file containing GFR namelists.
    %
  \item
    If neither of the above finds an input file, GFR will check the current run
    directory for the file \file{gfr.in} and use that file if it exists.
\end{enumerate}
%
If any of these steps finds what appears to be a valid GFR input file, it will
proceed with reading that input file to run a simulation.
%
If no input file is provided, the code will output the file
\file{gfr.in.default} which contains all of the available input
variables and their current default values.
 
%%%%%%%%%%%%%%%%%%%%%%%%%%%%%%%%%%%%%%%%%%%%%%%%%%%%%%%%%%%%%%%%%%%%%%%%%%%%%%%%

The contents of the user input file are described in
\sref{SEC:User_Input}.
%
The user may include every input variable in the input file.
%
However, the only \emph{required} entries are those inputs that
the user wants to have non-default values.

%%%%%%%%%%%%%%%%%%%%%%%%%%%%%%%%%%%%%%%%%%%%%%%%%%%%%%%%%%%%%%%%%%%%%%%%%%%%%%%%

Input variables may be added, removed, or renamed between versions.
%
However, the code will attempt to convert an older input file to the
current format.

%%%%%%%%%%%%%%%%%%%%%%%%%%%%%%%%%%%%%%%%%%%%%%%%%%%%%%%%%%%%%%%%%%%%%%%%%%%%%%%%

\devNotes{
The \inptext{nml\_version} namelist block is not required.
%
The \inp[nmlVersion::NamelistVersion]{Namelist\_Version} variable
defaults to the current value.
%
If the namelist file contains the \texttt{input} block,
the version is set to 1.
%
Otherwise, if the namelist file contains the \texttt{nml\_version} block,
it will attempt to read the value of
\inp[nmlVersion::NamelistVersion]{Namelist\_Version}
from that block.
}

%%%%%%%%%%%%%%%%%%%%%%%%%%%%%%%%%%%%%%%%%%%%%%%%%%%%%%%%%%%%%%%%%%%%%%%%%%%%%%%%

